\chapter{Introduction}
\label{ch:introduction}



\section{Overview}
\label{sec:intro:overview}


Randomized control trials represent a fast-growing and critical tool in empirical economics, boasting a literature that is vast, still growing \citep{Imbens2023}, and has recently birthed its own 
Nobel prizes \citep{Banerjee2020,Duflo2020,Kremer2020}. Much of the focus on this literature, however, is on the internal validity of the experimental design, as opposed to the external validity of results for policy relevance \citep{Heckman2008,Deaton2018,Manski2013}. This thesis concerns itself with two key components of external validity: representativity and ecological validity \citep{Neyman1934,Schmuckler2001}. It develops novel methodologies for leveraging the particular properties of digital ad platforms to address both of these components: (i) ecological validity when studying the impact of digital interventions and (ii) representativity more broadly, through optimizing participant recruitment subject to budget contraints. 

\section{External validity as a first-order concern}
\label{sec:intro:ev}

The relative paucity of focus on external validity has been noted by the field's own critics for decades. \citet{Manski2013} writes that ``analyses of experimental data have tended to be silent on the problem of extrapolating from the experiments performed to policies of interest'' (p.~37); the credibility of a study, in his account, depends jointly on internal and external validity (pp.~36--37). \citet{Heckman2008} echoes the diagnosis: forecasting ``the effects of interventions in environments other than the one where the estimate was obtained'' --- ``the problem of external validity'' (p.~8) --- is a task that randomised estimation, by construction, does not deliver. His structural account distinguishes between the statistical estimation problem (what RCTs solve) and the economic forecasting problem (what policy requires), and names the distance between them as the field's open question. \citet{Deaton2018} restates the critique in its strongest contemporary form: randomisation buys internal validity by design, but ``credibility in estimation can lead to incredibility in use'' (p.~3).

The recognition is not purely theoretical. A body of empirical work demonstrates that randomised estimates often fail to replicate across contexts, populations, and time periods (\cite{Pritchett2016, Allcott2015, Bisbee2017, Rosenzweig2019}). Additionally, the scale-up literature has documented the pilot-to-scale voltage drop empirically \citep{Bold2018} and programmatically \citep{Muralidharan2017,Al-Ubaydli2020}. Even from inside the randomised-controlled-trial movement itself, external validity is acknowledged as unfinished business. \citet{Snowberg2016} write that ``creating a rigorous framework for external validity is an important step in completing an ecosystem for social science field experiments.'' 




\section{Defining External Validity and Experiment Types}
\label{sec:intro:definitions}

\citet{CampbellStanley1963} decompose external validity across four sub-dimensions: populations, settings, treatment variables, and measurement variables. We will use the term \emph{representatvity} to address the dimension of ``population.'' In the standard sense from \citet{Neyman1934}, a sample is representative of a population if the sample's joint distribution, on relevant covariates, matches that of the target population of interest. We will use the term \emph{ecological validity} to bundle the remaining three sub-dimensions from \citet{CampbellStanley1963}: the setting, stimulus, and response to the conditions of an intervention. This is aligned with the three-dimensional definition of the term in \citet{Schmuckler2001}.

\citet{Harrison2004} similarly define six properties to determine the field context of an experiment, which can broadly map to representativity and ecological validity: (i) the subject pool, (ii) the prior knowledge of the subject pool, (iii) the item, (iv) the task, (v) the stakes, (vi) the environment. They also lay out a taxonomy for field experiments which will be useful for defining much of the work we will do in this thesis: [TODO: add field experiment types]. 

\section{Digital Recruitment}
\label{sec:intro:definitions}

TODO: We here need to make the point that people are using mturk and prolific to get study participants. This maps broadly onto the same sort of concerns that were originally used to raise the external validity concerns, that of using students for all studies (Harrison2004). We need some citations to show pepole are doing that. 

We then need to put that in context of both represenativity (mturk not that great, prolific or othe rpanel providers may provide panels that are representative across a fixed set of variables in certain countries, at high cost) and ecological validity (usually low, given that these individuals are generally asked to do something outside of their natural course of actions). 

Note that representativity is not a binary concept nor an end in and of itself. Then go into the basics of Ch1, where we define it as a continuous variable that can be optimized subject to budget constraints. Note that we start with the case of direct surveys and population estimation as the basic case. This is a methodological contribution in terms of an idea. 

It is only an interesting idea because we have a sequential problem. This is the recruiting via ad platforms. 

All this should be pulled from Ch 1 in some degree. 

\section{Experiments of Advertising for Public Good}

TODO: Here we need to cite the literature, including the covid ad lit, and whatever else is pulled from Ch2 lit review and intro. 

Importantly, we need to set up the gap. We then frame the gap in terms of ecological validity (mostly) or representativity (not sure if anything is). 

Then we introduce the two experiment design methods in the Ch2 paper and explain them in terms of ecological validity (feed experiment) and scale-up (cluster trial with holdout clusters).

\section{Experiments of Digital Apps for Public Health}

TODO: Similar quick lit citing here based on ch3 lit refs. Again, the main work here is to put it in context of representativity or ecological validity. 

Then we introduce the experiment within ch3 in terms of representativity (multi-country) and ecological validity (real parents without a previous interst in the app). 


%Digital experiments delivered through real platforms deliver on all three of these ecological-validity dimensions in ways no classical methodology matches. The \emph{setting} is the environment users actually inhabit --- their social-media feed, their app store, their messaging inbox. This is not a lab, not a survey frame, not a vignette. The \emph{treatment} is delivered through the same channel and form users encounter daily --- an ad in the feed, an app install prompt, a message. Not a lab stimulus or a printed handout. The \emph{response} is the behaviour the intervention targets, collected through the same infrastructure. This bundle has a name in experimental economics: \citet{Harrison2004} define a \emph{natural field experiment} as one in which subjects undertake naturally occurring tasks in their naturally occurring environment, unaware they are in an experiment. Their six-dimension field-context taxonomy --- subject pool, information, commodity, task, stakes, and environment --- gives precise language for which dimensions match the policy context and which do not. The ecological-validity gain over classical in-survey vignettes, lab analogs, onnr framed-field experiments is genuine and well-documented in the ad-as-treatment lineage \citep{Athey2023,Breza2021,Ho2023}.




\section{The platform tension: ecological validity and representativity as a joint problem}
\label{sec:intro:tension}

The unifying claim of this thesis is that digital platforms create a unique joint problem and opportunity for empirical economics: they simultaneously deliver an opportunity to affordably and quickly run studies with a high level of ecological validity, while systematically creating challenges to representativity. The two properties come from the \emph{same} infrastructure. This thesis develops and uses methodologies for solving the problems and gaining access to the opportunities inherent in the modern world. 

% No existing literature does so: the ad-as-treatment lineage emphasises the ecological-validity gains without engaging the representativity cost; the algorithmic-fairness literature documents the bias without engaging the ecological-validity opportunity.

\subsection{What digital platforms give us: ecological validity by construction}

\citet{CampbellStanley1963} decompose external validity across four sub-dimensions: populations, settings, treatment variables, and measurement variables (p.~5). By \emph{representativity} we mean the persons-dimension property that the sample's joint distribution on outcome-relevant covariates matches that of the target population --- the standard sense from \citet{Neyman1934}. By \emph{ecological validity} we mean the bundle of the remaining three sub-dimensions on the experimental-design side: the resemblance of setting, stimulus, and response to the conditions an intervention would meet were it actually deployed --- the three-dimensional formulation of \citet{Schmuckler2001}.

Digital experiments delivered through real platforms deliver on all three of these ecological-validity dimensions in ways no classical methodology matches. The \emph{setting} is the environment users actually inhabit --- their social-media feed, their app store, their messaging inbox. This is not a lab, not a survey frame, not a vignette. The \emph{treatment} is delivered through the same channel and form users encounter daily --- an ad in the feed, an app install prompt, a message. Not a lab stimulus or a printed handout. The \emph{response} is the behaviour the intervention targets, collected through the same infrastructure. This bundle has a name in experimental economics: \citet{Harrison2004} define a \emph{natural field experiment} as one in which subjects undertake naturally occurring tasks in their naturally occurring environment, unaware they are in an experiment. Their six-dimension field-context taxonomy --- subject pool, information, commodity, task, stakes, and environment --- gives precise language for which dimensions match the policy context and which do not. The ecological-validity gain over classical in-survey vignettes, lab analogs, or framed-field experiments is genuine and well-documented in the ad-as-treatment lineage \citep{Athey2023,Breza2021,Ho2023}.

\subsection{The catch: the same infrastructure breaks representativity}

But the platform infrastructure that delivers this ecological validity also systematically undermines representativity. Digital experiments inherit their samples through platform delivery, and platform delivery is optimised for engagement and cost, not for population representativity. Algorithmic delivery introduces systematic skew that is independent of the researcher's design choices.

\citet{Ali2019} find that Facebook's ad delivery produces significant demographic skew by race, gender, and socioeconomic status even when the advertiser uses neutral targeting criteria: the platform's engagement optimisation, not advertiser intent, drives the bias. \citet{Burtch2025} analyse over 181,000 platform A/B tests and 3,200 Lift tests on Meta's platform and find that audiences in treatment and control arms are systematically imbalanced --- and even Lift tests, while internally valid, face external validity concerns because the population eligible for ad exposure is algorithmically selected. \citet{Boegershausen2025} conduct a meta-review of 133 published platform studies and find that researchers systematically fail to report, account for, or even acknowledge algorithmic selection. \citet{Koenecke2023} show that cost-optimised delivery under-serves Spanish-speaking welfare-eligible populations on SNAP campaigns: the platform's objective function rewards cheaper-to-reach audiences, and the populations who most need the intervention are systematically under-reached. \citet{lambrecht2019algorithmic} and \citet{sapiezynski2022algorithms} document how engagement-maximising delivery tilts exposure toward users more likely to click and cheaper to reach, rather than toward those with higher expected welfare gains.

This is not passive self-selection --- the classic survey non-response problem. It is an \emph{algorithmic force} that actively channels treatment away from certain populations. The platform is not a neutral delivery channel; it is an optimising agent whose objective function is orthogonal to representativity.

\subsection{The tension is structural, not incidental}

The two properties --- ecological validity and representativity failure --- come from the same infrastructure. The platform can deliver the treatment in the real feed \emph{because} it controls delivery infrastructure. The platform skews who receives the treatment \emph{because} it optimises that delivery for engagement and cost. You cannot get one without the other.

This is a joint problem. Digital experiments on platforms sit at a point in external-validity space that no classical method occupies. They have higher ecological validity than any lab, vignette, or survey-embedded experiment. They have lower representativity than a well-designed probability survey. Both facts are true simultaneously.

This tension is \emph{not treated as a connected problem} in the existing literature. The ad-as-treatment lineage emphasises the ecological-validity gains without engaging the representativity cost. The algorithmic-fairness literature documents the bias without engaging the ecological-validity opportunity. The transportability and generalizability literature develops formal methods for reweighting trial samples to target populations, but assumes the trial sample is \emph{given} --- it does not address the problem of \emph{designing} the sample in the first place, under platform-side algorithmic constraints.

The thesis takes this joint problem as its object. The next section walks the two sub-literatures the thesis engages and shows how this tension manifests concretely.

\section{How the digital literature reflects the platform tension}
\label{sec:intro:lit}

\subsection{Ad-as-treatment: a trilemma on the outcome-measurement side}
\label{sec:intro:ad-trilemma}

A line of recent work runs randomised experiments \emph{through} ad-serving platforms, using paid ads as the treatment in public-health and behaviour-change contexts. Three canonical papers --- \citet{Athey2023}, \citet{Breza2021}, and \citet{Ho2023} --- score as natural field experiments on the treatment-delivery side: real channel, real subjects, real environment, subjects unaware \citep{Harrison2004}. But each faces a trilemma on the outcome-measurement side: each closes one external-validity gap and opens another, and the open slot that neither occupies is the one this thesis targets.

\citet{Athey2023} conduct a meta-analysis of 819 Meta Brand Lift Studies on COVID-vaccine beliefs, using single-item survey responses administered seconds after ad exposure. The gain is scale --- 2.1 billion person-ad pairs, tightly integrated with platform infrastructure. The cost is a construct-validity gap: the brand-lift survey measures immediate recall and stated attitudes, not durable behaviour. The paper acknowledges the bridge to vaccination behaviour is an unvalidated cross-sectional surrogate and the funnel pattern in their results shows brand-lift methodology is best-powered for low-construct-distance items (Knowledge) and null on behaviour-proximal items (Willingness, Effectiveness).

\citet{Breza2021} run a geo-cluster RCT with Facebook ads as the treatment and mobility/location traces as the outcome. The gain is real off-platform behaviour, real stakes, subjects unaware their mobility is the outcome --- the cleanest natural field experiment on the output side. The cost is the outcome space is restricted to what telemetry can see: travel, yes; vaccination uptake, bed-net use, knowledge, practices --- no.

\citet{Ho2023} replicate the Breza pattern with administrative vaccination records in a two-country follow-up after Omicron: researcher-assigned geo-clusters, platform-executed delivery, off-platform outcome via a parallel data stream. The null result confirms the principle: you get the outcome construct your parallel data stream supports.

The three sit at distinct points on the same gap. Neither Athey (brand-lift) nor Breza/Ho (telemetry/records) can measure arbitrary behavioural outcomes --- bed-net use, care-seeking, knowledge, practice change --- at scale, on the population the campaign already acts on, through the same channel. That slot is open.

All three inherit the same representativity problem: the sample is whatever the platform's targeting, bidding, and engagement optimisation produce. The algorithmic targeting literature demonstrates that this inheritance carries systematic bias \citep{Ali2019,Boegershausen2025,Koenecke2023}. In the malaria-prevention campaign of Chapter~\ref{ch:malaria}, this manifested concretely: engagement-optimised delivery concentrated ad exposure among higher-SES, urban users --- precisely the populations with lower baseline malaria risk. The campaign's ecological validity was genuine; its representativity was not. The design could not address one failure without compromising the other --- a structural constraint, not a design oversight.

\subsection{App-based interventions: the existing-user constraint as a representativity failure}
\label{sec:intro:apps}

A sparser literature evaluates mobile app interventions in parenting and early-childhood development. \citet{Jaggi2025} find only five experimental studies in a scoping review of digital early-childhood parenting interventions in LMICs. Two systematic reviews identify self-selection and retention as critical weaknesses across the mHealth evaluation literature \citep{Sapanel2023,Bremer2023}.

What this literature delivers is genuine: an app is the real product in the real environment --- real subjects, real task, real channel --- a framed field experiment in the Harrison--List sense \citep{Harrison2004}. Ecological validity on the treatment-delivery dimension is genuine.

What it characteristically lacks is representativity. Four gaps recur with enough regularity to read as binding constraints:

\begin{enumerate}[leftmargin=*]
    \item \textbf{Samples constrained to existing users.} Classical app evaluations sample within the installed-user population. Non-users are excluded by construction. Health app users are younger, wealthier, and more educated than non-users \citep{Carroll}, and the full causal chain --- encouragement, install, engagement, outcomes --- cannot be observed. The binding policy question --- would non-users use this if encouraged? --- is unanswerable.
    \item \textbf{Single-country designs.} Most evaluations stay within one country and one app version, leaving cross-context generalisation undone. A scoping review of digital parenting interventions in LMICs identifies single-country designs as near-universal \citep{Jaggi2025}.
    \item \textbf{Short-horizon outcomes.} Knowledge and confidence are measured days to weeks after install; long-run behaviour change is rarely captured.
    \item \textbf{Attrition compounds selection.} Median retention in digital health studies is 5.5 days; those who stay differ systematically from those who drop \citep{Pratap2019}. The analytic sample is doubly selected --- once by adoption and again by sustained engagement.
\end{enumerate}

These gaps are structural, not incidental. Recruiting non-users at scale and randomising encouragement requires an infrastructure --- digital ad recruitment with representativity targeting --- that the app-evaluation literature typically does not possess. \citet{Bradlow1998} formalises the encouragement design as the solution to self-selected samples, but the method has almost never been applied to mobile app evaluation at multi-country scale with non-users in the sampling frame.

\subsection{Recurring patterns and the pilot-to-scale gap}
\label{sec:intro:patterns}

Across both sub-literatures, three gap patterns recur with enough regularity to read them as binding open problems:

\begin{enumerate}[leftmargin=*]
    \item \textbf{Representativity is not actively managed.} Both ad-as-treatment and app-evaluation studies inherit samples rather than designing them. The algorithmic targeting literature shows that \emph{inheriting} a platform sample means inheriting a systematic bias \citep{Ali2019,Boegershausen2025}.
    \item \textbf{Ecological validity is asserted, not designed.} Both literatures \emph{can} deliver ecological-validity properties --- real channel, real subjects, real task --- but most studies inherit whatever the platform provides. The Harrison--List vocabulary makes this distinction legible: it separates ``it's digital so it's real'' from ``the experimental architecture \emph{guarantees} a specific external-validity property.''
    \item \textbf{Pilot-to-scale voltage drop} is documented in the broader scale-up literature \citep{Muralidharan2017,Bold2018,Al-Ubaydli2020} but rarely engaged in digital-experimentation work. Digital papers often claim at-scale delivery implicitly --- the platform \emph{is} the scale --- without the experimental architecture to back the claim.
\end{enumerate}

A subtler running thread concerns \emph{demand characteristics}: subjects knowing they are in an experiment and modifying their behaviour. \citet{Corneille2023} review sixty years of demand-characteristics research and conclude that demand effects threaten both internal and external validity. \citet{deQuidt2018} find measurable demand effects even in anonymous online experiments, averaging 0.13 standard deviations across 11 classic experimental tasks with approximately 19,000 participants. Most digital-experiment studies fall on the wrong side of this line: subjects consent, subjects know they are being studied, demand effects are uncontrolled. The Harrison--List distinction between framed and natural field experiments is not taxonomic pedantry --- it marks an external-validity boundary that most of the literature crosses without comment.

These recurring patterns define the contribution space. The thesis develops methodologies that treat representativity as a controllable design parameter and ecological validity as a property \emph{guaranteed by experimental architecture} rather than inherited from the platform.

\section{Three contributions}
\label{sec:intro:contributions}

Each chapter addresses a distinct piece of the joint problem identified in Section~\ref{sec:intro:tension}. The chapters are connected not merely by shared digital infrastructure but by their relationship to the platform tension: Chapter~\ref{ch:survey-sampling} solves the representativity problem that the platform creates; Chapter~\ref{ch:malaria} solves the ecological-validity problem that the platform partially delivers and partially undermines; Chapter~\ref{ch:bebbo} demonstrates that, with the right recruitment infrastructure, a classical experimental design can overcome the existing-user constraint that limits the app-evaluation literature. The chapters are related in substance, not just by adjacency.

\subsection{Representativity: a unified adaptive recruitment framework}
\label{sec:intro:ch1}

Chapter~\ref{ch:survey-sampling} poses representativity as a \emph{budget-constrained optimisation problem} that operates \emph{during} recruitment, not after it. Three operational capabilities distinguish the approach:

\begin{enumerate}[leftmargin=*]
    \item \textbf{Active chasing of under-represented strata during recruitment.} Traditional surveys recruit and then weight what they got; this procedure allocates while recruiting, dynamically responding to the stratum-level composition of the incoming sample.
    \item \textbf{Sequential variance estimation per stratum,} updated as data arrives, driving allocation --- a dynamic analogue of Neyman optimal allocation \citep{Neyman1934}.
    \item \textbf{Recruitment and estimation fused into one optimisation} --- classically separate steps, unified here. Allocation decisions are informed by the same estimation objective the final analysis targets.
\end{enumerate}

The procedure anchors in the adaptive experimental design literature \citep{Kasy2021,Caria2020,Carneiro2020}. The binding use case is digital recruitment through Meta's ad platform, where cost and strata-level response rates are estimated dynamically and allocation adjusts in real time. But the formulation is generic: surveys are the simplest demonstration; experiment recruitment --- what Chapter~\ref{ch:malaria} and Chapter~\ref{ch:bebbo} depend on operationally --- is the natural extension.

This is an external-validity contribution because the procedure makes the sample's joint distribution of observables match the target population --- representativity as the \emph{design constraint}, not an afterthought. Critically, it operates \emph{against} the platform's algorithmic selection bias: the platform optimises for engagement; the recruitment procedure optimises for representativity. This is the only way to get digital ecological validity without paying the full representativity cost.

\subsection{Ecological validity: the feed experiment}
\label{sec:intro:ch2}

Chapter~\ref{ch:malaria}'s methodological contribution is the \emph{feed experiment}: using the saturated ad surface as the survey-administration channel to recover arbitrary behavioural outcomes affordably from the population the campaign already acts on. This is the open slot in the trilemma identified in Section~\ref{sec:intro:ad-trilemma}. Neither \citet{Athey2023} (on-platform brand-lift) nor \citet{Breza2021}/\citet{Ho2023} (off-platform telemetry/records) occupies it.

The design has two properties. First, it is a natural field experiment on both delivery \emph{and} measurement through the same channel. The treatment is delivered through the same feed the campaign saturates; the outcome survey is administered through the same feed --- no separate survey-operations apparatus, no off-platform follow-up. Subjects are unaware of the experiment at both delivery and measurement stages. This is more natural-field on the output side than \citet{Athey2023} achieve (brand-lift survey is an explicit instrument subjects consciously interact with) without acquiring the restriction to whatever outcome construct telemetry can see that limits \citet{Breza2021}. Second, the feed survey can measure arbitrary behavioural outcomes: bed-net use, care-seeking intentions, knowledge change, practice change --- any construct a household survey can ask about. This generality is what brand-lift cannot provide and telemetry cannot reach.

The underlying architecture is a cluster-randomised design carved from the actual national rollout --- clusters sample the policy as it operates at scale \citep{Hayes2017}. The cluster-RCT architecture follows the Breza lineage (researcher-side cluster assignment; platform executes delivery within pre-assigned zones). The novelty is the feed-experiment measurement move on top.

The integration with Chapter~\ref{ch:survey-sampling} is tighter than shared infrastructure. Chapter~\ref{ch:survey-sampling}'s recruitment machinery is what makes the feed-experiment surveys land on a meaningful sample inside the saturated population. The campaign creates the saturated population; Chapter~\ref{ch:survey-sampling} ensures the survey sample within it is representative rather than convenience-drawn. Chapter~\ref{ch:survey-sampling} solves the representativity problem \emph{for} Chapter~\ref{ch:malaria}'s ecological-validity design, and neither chapter's contribution works without the other.

Substantively, the campaign raised reported bed-net use from 69\% to 74\% on average, with the cluster effect concentrated in pucca (durable-material) dwellings --- reach, not persuasion, was the binding constraint. The SES-differentiated effect is evidence of the platform tension in action: ecological validity of delivery was genuine; representativity of exposure failed. The feed-experiment methodology disentangles these two.

\subsection{Feasibility application: encouragement at scale}
\label{sec:intro:ch3}

Chapter~\ref{ch:bebbo} applies a classical design to a problem that the literature has not reached. The design is standard: an Imbens--Angrist--Rubin encouragement design with two-sided non-compliance, ITT and ToT effects recovered under the standard monotonicity and exclusion assumptions \citep{Imbens2015}. The contribution is \emph{feasibility enabled by the digital recruitment infrastructure} that Chapter~\ref{ch:survey-sampling} develops.

Three features distinguish the design from the existing app-evaluation literature. First, \emph{multi-country scale}: caregivers are recruited across multiple countries through Meta's ad platform, in a single coordinated operation. Second, \emph{non-users in the sampling frame}: classic app evaluations are constrained to existing users; encouragement at scale lets the study include non-users and capture the full causal chain --- encouragement, install, engagement, outcomes \citep{Bradlow1998}. Third, \emph{medium-term follow-up}: Chapters~\ref{ch:malaria}'s feed-experiment structure collects unconsented outcomes at the population level through the same ad channel; Chapter~\ref{ch:bebbo} collects consented longitudinal individual-level outcomes that the natural-FE posture of Chapter~\ref{ch:malaria} cannot. The trade-off is explicit: Chapter~\ref{ch:bebbo} is a framed field experiment in the Harrison--List sense (real subjects, real task, real channel, but subjects consent and know they are in a study), and in return it collects rich individual-level panel data.

Substantively, the null is on parenting knowledge, confidence, and self-reported practices; 76\% day-one abandonment points to an engagement bottleneck, not a content failure. The contribution is space-specific: \citet{Jaggi2025} identify only five experimental studies of digital parenting interventions in LMICs; none combine encouragement design, multi-country scale, and non-users in the sampling frame.

\section{Roadmap}
\label{sec:intro:roadmap}

The remainder of the thesis is organised as follows.

Chapter~\ref{ch:survey-sampling} develops the Theme-1 contribution: an adaptive stratified-allocation procedure with post-stratification estimation, simulation evidence, and an applied case study of survey recruitment through Meta's ad platform.

Chapter~\ref{ch:malaria} develops the feed experiment: a district-cluster-randomised behaviour-change campaign on bed-net use in rural India, with panel outcome waves administered through the same ad channel. Cluster-level treatment effects are estimated under the standard cluster-RCT framework.

Chapter~\ref{ch:bebbo} executes the encouragement design: multi-country recruitment of caregivers via Meta, random encouragement to install the Bebbo parenting app, and ITT/ToT estimation of effects on parenting knowledge and self-reported practices, attending to attrition and two-sided non-compliance.

Chapter~\ref{ch:conclusion} returns to the platform tension, draws the cross-paper synthesis, and discusses limitations and the design space that the feed-experiment idea opens for future work.

\printbibliography[heading=subbibliography, title={References for Chapter \thechapter}]
